\documentclass[a4paper,12pt]{article}
\usepackage[UTF8,fontset=none]{ctex}
\usepackage[top=1.5cm,bottom=1.5cm,left=2cm,right=2cm]{geometry}
\usepackage{array}
\usepackage{makecell}
\usepackage{fontspec}

\setmainfont{TeX Gyre Termes}
\setCJKmainfont{Source Han Serif CN}
\setCJKsansfont{Noto Sans CJK SC}

\pagestyle{empty}
\setlength{\parindent}{0pt}
\renewcommand{\arraystretch}{1.4}

\begin{document}

\begin{center}
{\Large\bfseries 湖北文理学院毕业设计（论文）成绩评定表}
\end{center}

\vspace{6pt}

\begin{tabular}{|p{1.2cm}|p{2.2cm}|p{1.2cm}|p{2.5cm}|p{1.5cm}|p{5.06cm}|}
\hline
姓名 & 廖嘉旺 & 学号 & 2022117117 & 专业班 & 计科2211 \\
\hline
题目 & \multicolumn{5}{l|}{基于Apollo的多障碍物场景下路径规划算法实现} \\
\hline
\makecell{指\\导\\教\\师\\评\\语} & \multicolumn{5}{p{13.46cm}|}{
\vspace{4pt}
课题主攻方向明确，研究路线切实可行，技术手段选择合理，开题报告格式规范；能很好地完成任务书规定的工作量，学习态度认真；按时完成外文翻译，译文质量高；对研究的问题能较深刻分析或有独到之处，成果应用价值高，反映出很好地掌握了有关基础理论与专业知识；论文结构严谨，逻辑性强，论述层次清晰，语言准确，文字流畅，论文符合规范要求。
\vspace{8pt}

得分：\hspace{2cm} 签字：\hfill 2025年5月7日\hspace{0.5cm}
\vspace{2pt}
} \\
\hline
\makecell{评\\阅\\教\\师\\评\\语} & \multicolumn{5}{p{13.46cm}|}{
\vspace{40pt}

得分：\hspace{2cm} 签字：\hfill 2025年5月7日\hspace{0.5cm}
\vspace{2pt}
} \\
\hline
\makecell{答\\辩\\小\\组\\评\\语} & \multicolumn{5}{p{13.46cm}|}{
\vspace{40pt}

得分：\hspace{2cm} 组长签字：\hfill 2025年5月23日\hspace{0.5cm}
\vspace{2pt}
} \\
\hline
\makecell{学\\院\\意\\见} & \multicolumn{5}{p{13.46cm}|}{
\vspace{2pt}
经学院答辩委员会综合评定，该同学毕业论文（设计）成绩为：
\vspace{12pt}

\hfill 学院（公章）\hspace{2cm}

\hfill 2025年5月23日\hspace{0.5cm}
\vspace{2pt}
} \\
\hline
\end{tabular}

\end{document}
